\section{Discussion, External Validity, and Limitations}
\label{sec:discussion}

The empirical findings from our primary frozen benchmark ($H_1$--$H_4$) and cross-schema study demonstrate that learned adaptive schema resolution with reactive recovery provides an efficient, reliable foundation for enterprise tool agents. Below, we analyze architectural boundaries, validity threats, and operational trade-offs.

\subsection{Architectural Scope and External Validity}
\noindent\textbf{1) Schema Topologies and Structural Expressiveness}: Enterprise APIs exhibit diverse topologies; our evaluation incorporated nested objects and arrays (e.g., line items in \texttt{create\_purchase\_order}, packages in \texttt{shipment\_creation}). While JSON Pointer traversal handles nested definitions deterministically, polymorphic branches (\texttt{oneOf}, \texttt{anyOf}) introduce pruning risks if discriminators are omitted. In LASR, discriminator fields and structural keys are designated invariants ($C_I^*$), keeping schema branching unambiguous. Automated DAG analysis for recursive schemas remains future work.

\noindent\textbf{2) Schema Validator Standards vs. Business Logic}: A core design philosophy of LASR-MCP is strict adherence to JSON Schema Draft 2020-12 specifications~\cite{jsonschema202012}. The validator reports error locations as JSON Pointer paths~\cite{rfc6901} and validation keywords; LASR maps these diagnostics to schema-specific canonical constraint nodes without bespoke heuristics. Enterprise backends enforce transactional business rules (e.g., stock levels, credit checks) outside schema syntax; LASR-MCP explicitly scopes itself to syntactic and structural schema conformance, leaving business failures to application compensation.

\noindent\textbf{3) Model-Family Sensitivities and Parameter Scale}: Controlled experiments benchmarked Google's \texttt{gemini-3.5-flash-lite} as a representative enterprise model. While pilot ablations observed unconstrained enum generation across evaluated configurations, this empirical observation does not establish that the phenomenon is model-independent, and we make no validated transfer claims to other model families or parameter scales. Conceptually, lean schemas ($C_I^*$) reduce attention dilution in smaller models while mitigating multi-turn token costs in frontier models; validating cross-family transfer remains an important empirical objective.

\noindent\textbf{4) Unit of Analysis and Constraint Recurrence}: Our primary $H_1$ prediction evaluation ($N=32$) reflects paired task-level accuracy across held-out template family C. As disclosed in Section~\ref{sec:results}, all four target constraint combinations recur across train and test splits, and vocabulary overlap is 69.1\% due to domain-specific enterprise terminology. A four-cluster sign-flip sensitivity analysis yields $p = 0.125$. Consequently, our findings support lexical predictive accuracy for this held-out template family with recurring constraint bundles, but do not claim inductive generalization to completely novel semantic constraint combinations.

\subsection{Empirical Provenance and Sensitivity Replications}
\noindent\textbf{1) Epistemological Separation and Benchmark Provenance}: We emphasize strict separation of empirical scopes: \textbf{Scope 1 (Primary Frozen Benchmark, $N=32$ Tasks)}: Evaluates $H_1$--$H_4$ on the frozen procurement benchmark (commit \texttt{9807340}, SHA-256 \texttt{78c2c88e}) under live Google \texttt{gemini-3.5-flash-lite} execution, establishing predictive generalization ($H_1$), invariant safety ($H_2$), recovery lift ($H_3$), and regime-dependent token efficiency ($H_4$). \textbf{Scope 2 (Post-Freeze Mechanism Stress Test, $N=32$ Tasks, 64 Executions)}: Evaluates algorithm and architecture transfer across three new synthetic schemas plus primary replication (four schemas total, SHA-256 \texttt{b3d57636}) under a deterministic simulation-based stress test with heuristic token counts, providing preliminary evidence of recovery mechanism robustness without live model API billing or altering the primary benchmark. \textbf{Scope 3 (Canonical Live Benchmark, $N=3,840$ Trials, 7,860 Calls)}: Evaluates full-scale systems performance across 32 tasks, 6 arms, and 20 fixed-seed repetitions on \texttt{gemini-3.5-flash-lite} at temperature $0.2$, comprising 7,860 response-bearing calls and 154 client transport retries (archive digest \texttt{fc82d33b\allowbreak 8992914e\allowbreak e130ef28\allowbreak 05427ef2\allowbreak a4f01541\allowbreak 243e0f93\allowbreak be120e2d\allowbreak 20ae7f5b}).

\noindent\textbf{2) Exploratory Live Replication and Sensitivity Analysis}: Prior to executing the canonical $N=3,840$ live benchmark, we conducted an exploratory single-pass replication on the Google Generative Language API using \texttt{gemini-3.5-flash-lite} across all 32 tasks and 6 arms to assess endpoint stability. This preliminary run is explicitly distinguished from the definitive $N=3,840$ dataset by three measurement bounds: (1) recorded executions produced 428 attempt records (285 newly observed unique live response entries written on initial misses alongside approximately 143 within-run prompt-cache reuses); (2) turnaround latencies reflect client logical wall time (including HTTP retries, backoff, and pacing sleeps); and (3) requests did not transmit \texttt{seed=42}, and cache records lack provider model-version strings or full \texttt{usageMetadata}. Under this exploratory run, observed prompt-plus-candidate token subtotals were: Full Schema 55,549 tokens; CI $k=0$ 69,555 ($-25.21\%$); CI $k=1$ static 59,445 ($-7.01\%$); LASR $k=1$ reactive 52,389 ($+5.69\%$); LASR $k=2$ reactive 51,224 ($+7.79\%$); and Oracle reactive 53,415 ($+3.84\%$). Conditional mean structural $TTS_{\text{struct}}$ tokens among structurally successful tasks were: Full 1,694.7; CI $k=0$: 2,142.5 ($-26.42\%$); CI $k=1$ static: 1,738.0 ($-2.56\%$); LASR $k=1$ reactive: 1,603.7 ($+5.37\%$); LASR $k=2$ reactive: 1,534.8 ($+9.43\%$); and Oracle: 1,617.0 ($+4.58\%$). First-pass hierarchical success ($\text{structural} \land \text{semantic}$) was: Full 13/32; CI $k=0$ 0/32; CI $k=1$ static 2/32; LASR $k=1$ reactive 2/32; LASR $k=2$ reactive 4/32; and Oracle 2/32. Paired bootstrap CIs against Full cross zero in this sample (LASR $k=1$~$[-13.2\%, +24.4\%]$; LASR $k=2$~$[-11.0\%, +26.3\%]$), reflecting hosted model variability. We treat this rerun strictly as an exploratory sensitivity replication rather than confusing it with canonical $N=3,840$ findings.

\subsection{Break-Even Boundary and Analytical Scaling Projections}
\noindent\textbf{1) The Retry-Amortization Break-Even Boundary}: The full-scale $N=3,840$ live benchmark (32 tasks, 6 arms, 20 fixed-seed executions per task-arm pair at temperature $0.2$) characterizes an observed \emph{Retry-Amortization Boundary} across the four evaluated enterprise task archetypes in agentic tool-calling. Let $P_{\text{base}}$ denote base user prompt tokens, $S_{\text{full}}$ monolithic schema tokens, $S_0$ upfront lean schema tokens, and $\Delta S$ reactive constraint expansion. Let $C_{\text{gen},1}$ and $C_{\text{gen},2}$ denote candidate generation lengths across turns (with $C_{\text{gen},1} \approx C_{\text{gen},2} \approx C_{\text{gen}}$).
In single-turn success ($T=1$), token savings is strictly positive ($\Delta \text{Tok} = S_{\text{full}} - S_0 > 0$). When an omitted constraint forces a validation error and retry ($T=2$), the request re-transmits base prompt $P_{\text{base}}$, expanded schema $S_0 + \Delta S$, prior history $H_1$ ($|H_1| \approx C_{\text{gen}}$), and validator feedback $F_1$, yielding modeled two-turn cost:
\begin{equation}
\label{eq:cost_t2}
\text{Cost}(T=2) \approx 2 P_{\text{base}} + 2 S_0 + \Delta S + H_1 + F_1 + 2 C_{\text{gen}}
\end{equation}
which exceeds the single-turn monolithic cost $\text{Cost}_{\text{full}}(T=1) = P_{\text{base}} + S_{\text{full}} + C_{\text{gen}}$ whenever:
\begin{equation}
P_{\text{base}} + S_0 + \Delta S + H_1 + F_1 + C_{\text{gen}} > S_{\text{full}} - S_0
\end{equation}
Terms in this derivation represent modeled analytical estimates. When retry turn overhead $(P_{\text{base}} + S_0 + \Delta S + H_1 + F_1 + C_{\text{gen}})$ exceeds the upfront compression delta $(S_{\text{full}} - S_0)$ as formalized in Inequality~(15), a multi-turn retry cascade eliminates the initial token savings. Note that a simpler reactive baseline that falls back to the full monolithic schema ($S_{\text{full}} = 552.0\text{ tok}$) on Attempt 2 rather than splicing only the violated delta ($S_0 + \Delta S = 404.6\text{ tok}$) would incur an additional $+147.4$ schema tokens per retry turn.

Empirically across our 32 task clusters, this boundary creates a sharp domain divergence: (1)~\textbf{Routine Regimes ($C_1, m_{\text{task}}(x) \le k$)}: In Cluster 1 (\texttt{routine\_regional\_replenish}), active task constraints fall within the $k=2$ upfront bundle ($75.0\%$ first-pass success). Initial schema compression ($359.4$ vs.\ $552.0$ tokens) is preserved without round-trip penalties, delivering approximately $\mathbf{13.5\%}$ lower $P+C$ token usage ($847.6$ vs.\ $979.8$ mean recorded tokens) at $100.0\%$ structural reliability (signed savings: $+13.5\%$). We do not generalize this cluster-specific result to the full benchmark. (2)~\textbf{Deep Inter-Field Regimes ($C_2$--$C_4, m_{\text{task}}(x) > k$)}: In complex procurement tasks with cross-field dependencies, omitted constraints trigger Attempt 1 failures ($0.0\%$ first-pass). Each retry re-transmits prompt and prior generation, incurring multi-turn penalties that inflate net workload tokens by $+22.2\%$ ($C_2$, $1,557.5$ vs.\ $1,274.2$; signed savings: $-22.2\%$), $+59.8\%$ ($C_3$, $1,610.9$ vs.\ $1,008.1$; signed savings: $-59.8\%$), and $+33.6\%$ ($C_4$, $1,609.9$ vs.\ $1,205.0$; signed savings: $-33.6\%$) relative to Full Control. Across the full 32-task benchmark ($N=3,840$ trials), Full Control achieves a mean recorded subtotal of $1,116.8$ tokens versus $1,406.5$ for LASR $k=2$ (workload increase: $+25.94\%$, $95\%$ CI: $[+3.16\%, +52.55\%]$; signed relative savings: $-25.94\%$, $95\%$ CI: $[-52.55\%, -3.16\%]$, bootstrap CI excluding zero); we do not describe the overall benchmark result as token savings.

\noindent\textbf{2) Multi-Tool Catalog Analytical Projections}: We emphasize that our empirical evaluation evaluated single-tool isolation with deep constraint hierarchies. In multi-tool agent registries where an agent is exposed to $W > 1$ heterogeneous enterprise tools $\mathcal{T} = \{t_1, \dots, t_W\}$ concurrently (e.g., $W = 10$ or $W = 50$ enterprise tools in an MCP catalog), each tool $i$ possesses monolithic schema size $S_{\text{full},i}$ and pruned lean schema size $S_{0,i}$. The gross upfront catalog savings from compressing all $W$ tools simultaneously and net savings under single-tool invocation failure are modeled as:
\begin{align}
\label{eq:multitool_savings}
\text{Savings}_{\text{catalog}}^{\text{gross}} &= \sum\nolimits_{i=1}^W \bigl( S_{\text{full},i} - S_{0,i} \bigr), \\
\label{eq:multitool_net}
\text{Savings}_{\text{net}} &= \text{Savings}_{\text{catalog}}^{\text{gross}} - \mathbb{P}(\text{retry}) \cdot \text{Penalty}_{\text{retry}}.
\end{align}
Under homogeneous average tool dimensions ($S_{\text{full},i} \approx \bar{S}_{\text{full}}$, $S_{0,i} \approx \bar{S}_0$), this scales as $W \cdot (\bar{S}_{\text{full}} - \bar{S}_0)$. Net savings remains strictly positive whenever $\sum_{i=1}^W (S_{\text{full},i} - S_{0,i}) > \mathbb{P}(\text{retry}) \cdot \text{Penalty}_{\text{retry}}$, where $\text{Penalty}_{\text{retry}} \approx P_{\text{base}} + S_0 + \Delta S + H_1 + F_1 + C_{\text{gen}}$. Because gross upfront savings scales linearly with $W$ while the retry penalty applies only to the single invoked tool, multi-tool registries are analytically projected to yield substantial net savings even under modest retry rates. We restrict this scaling relationship strictly to an analytical projection and future research direction, while keeping our formal empirical claims strictly grounded in the single-tool isolation dataset ($N=3,840$).

\subsection{Limitations \& Methodological Disclosures}
We explicitly disclose four methodological boundaries of this study: (1)~\textbf{Fixed-Seed Determinism}: All 20 repetitions per task-arm combination produced identical first responses and outcomes under the repeated fixed seed (temperature $0.2$, seed 42), verifying execution stability rather than representing independent stochastic draws. (2)~\textbf{Data Collection Timeline}: As detailed in Section~\ref{sec:setup}-D, collection spanned multiple days and sessions across paid and quota-limited API accounts (completing 3,149 initial trials followed by a 2.8-hour final session completing the remaining 691 trials across backup keys). (3)~\textbf{Single-Tool Scope}: Empirical metrics reflect single-tool isolation rather than concurrent multi-tool exposure. (4)~\textbf{Model Specificity}: Live evaluations reflect Google's \texttt{gemini-3.5-flash-lite}; cross-family transfer remains subject to future work.

\section{Conclusion}
\label{sec:conclusion}
This study resolves an architectural question in agentic tool calling: when does lazy schema compression save tokens, and when do reactive retry cascades outweigh upfront reduction? Across $3,840$ completed empirical trials (32 tasks, 6 arms, 20 fixed-seed executions at temperature $0.2$), naive constraint stripping (\texttt{ci\_k0}) proved an operational anti-pattern that inflates net token consumption by $84.3\%$ while leaving $12.5\%$ of workflows structurally broken. By introducing prioritized $k=2$ upfront constraint bundling coupled with monotonic reactive subtree slicing (\texttt{ci\_k2\_reactive}), LASR-MCP cuts token waste by $31.7\%$ over naive lean baselines, restores $96.9\%$ structural completion, and operates within $2.1\%$ of a task-aware oracle baseline ($2.08\%$ higher relative to oracle). Crucially, our break-even analysis identifies an operational boundary governed by the ratio of task-conditioned constraints $m_{\text{task}}(x)$ to upfront budget $k$: in routine regimes ($C_1, m_{\text{task}}(x) \le k$), LASR $k=2$ delivers approximately $13.5\%$ lower $P+C$ token usage ($847.6$ vs.\ $979.8$ mean tokens) at $100.0\%$ reliability. Conversely, across the full 32-task benchmark, Full Control achieves a mean recorded subtotal of $1,116.8$ tokens versus $1,406.5$ for LASR $k=2$ (workload increase: $+25.94\%$, $95\%$ CI: $[+3.16\%, +52.55\%]$; signed relative savings: $-25.94\%$, $95\%$ CI: $[-52.55\%, -3.16\%]$, bootstrap CI excluding zero). Monolithic schemas thus remain optimal for isolated single-tool calls with deep cross-field invariants ($C_2\text{--}C_4$, $m_{\text{task}}(x) > k$) under single-tool objectives by avoiding retry cascades unless amortized across multi-tool registries. These empirical boundaries equip enterprise MCP architects with concrete, mathematically grounded deployment rules.

\section*{Declaration of Generative AI and AI-Assisted Technologies}
In accordance with IEEE and academic publishing guidelines on Generative AI transparency, the author discloses: (1) As the primary empirical subject, live API calls were executed against Google's \texttt{gemini-3.5-flash-lite} model ($N=3,840$ trials, $7,860$ API invocations) to evaluate tool-calling schema resolution under the \texttt{mcp-inventory-bench} v2.0.3 protocol. (2) During manuscript preparation, AI-assisted coding and editing tools were utilized strictly for LaTeX typesetting, copy-editing for readability, and generating deterministic plotting scripts (\texttt{TikZ}/\texttt{PGFPlots} and \texttt{matplotlib}) directly from verified CSV/JSON telemetry. No generative AI image-synthesis tools were used to create or alter figures. The author reviewed, audited, and cryptographically verified all code, empirical data, statistical bootstrap claims, and visual representations, assuming full responsibility for manuscript integrity.



